\section{Conclusion}
\label{sec:conclusion}

In \casepath, a model reads each deciding fact several times and quotes the file; code admits agreeing values, executes the organisation's rules and holds what the file leaves open. In a pre-registered test of a frozen commit on synthetic claims, the design held what single prompts with the same rules did not hold. On the cost-efficient model it had a lower mean error than every single-prompt and prior-work design in both data sets. The registered tests confirmed this against the single prompt in both and against two of three prior-work strategies on theft; against self-consistency the theft lead (\sthreeHFourSCA{} against \sthreeHFourSCB) gave $p\sthreeHFourSCPeq$, and that hypothesis failed. On the frontier model, run once, it made fewer tenancy errors and had no advantage on theft (\sthreeSecSevenAA{} errors per pair against \sthreeSecSevenAB{} and, with the original rules, \sthreeRicOrigFrontierThErrors), where the rules set the limit. The theft rules had been reviewed by AI sessions that knew the test conditions, and the tenancy references come from the engine \casepath implements; under one rejected reading of the rule text the product's unchanged workflow loses its tenancy lead. Whether it helps case handlers on real files is the question of the deployment study.
